% ======================================================================
% SECTION 2: PATIENT PRIORITY FRAMEWORK
% Location in paper: Second body section, sets the operational
% definition of patient control that all seven bills will reference.
% ======================================================================

[PROCESSING INSTRUCTION (Patient Priority Framework, 5 to 7 paragraphs
plus one ASCII diagram): Establish the operational definition of
"patient control" that every bill in this paper builds on. The
definition has seven functional dimensions, the first six adapted from
patients/paper/Deep-Research-3/part1_legal_baseline.md and the seventh
introduced new in this paper.

Read the following in full and synthesize:
- patients/paper/Deep-Research-1/chunk_1_foundations.md
  (the four-part taxonomy of patient control: informational,
  behavioral, relational, decisional; CASE-cancer, PAM, and self-
  efficacy validated instruments)
- patients/paper/Deep-Research-1/chunk_2_advanced_systems.md
  (QOL+, PATH/MyCoach, the software-only baseline against which any
  next-generation system must be evaluated; the ceiling-effect finding
  from PERSIST)
- patients/paper/Deep-Research-3/part1_legal_baseline.md (the six
  functional dimensions)
- patients/paper/Deep-Research-3/part3_metrics_guardrails.md (the
  four-domain dashboard: access, decision quality, control-in-action,
  fairness and technical quality)

Cover, in order, with one paragraph per dimension:

(2.1) Trial discovery: machine-readable eligibility, AI matching, no
provider gatekeeping. Cite \cite{jin2024trialgpt}, \cite{gupta2024prism},
and \cite{abdallah2026trialmatchai}.

(2.2) Enrollment comprehension: key information rules, broad consent,
patient-initiated translation. Cite \cite{fda_informed_consent_guidance_2023},
\cite{fda_ohrp_key_information_draft_2024}, and
\cite{ohrp_broad_consent_attachment_c_2017}.

(2.3) Location and scheduling choice: 24/7 booking, decentralized
elements, home-based delivery. Cite \cite{fda_dct_guidance_2024},
\cite{fda_oce_advancing_oncology_dct_2024},
\cite{dronca2026cancer_care_beyond_walls}, and
\cite{kawchak_2026_19994945} (for the prior 24/7 trial paper).

(2.4) Data and biospecimen governance: HIPAA, broad consent,
information-blocking penalties. Cite \cite{hhs_right_to_access_2025},
\cite{onc_cures_act_final_rule_page}, and
\cite{federal_register_info_blocking_disincentives_2024}.

(2.5) Toxicity warnings and adverse event self-management: ePRO,
ASyMS, eRAPID. Cite \cite{rocque2025remote_symptom_monitoring},
\cite{kearney2009asyms}, \cite{maguire2017esmart}, and
\cite{absolom2021erapid}.

(2.6) AI-exclusion contestability: HTI-1 DSI transparency, FDA AI
draft guidance 2025, mandatory appeal rights. Cite
\cite{rank9_onc_hti1_2023}, \cite{hti1_dsi_fact_sheet_2023}, and
\cite{fda_ai_regulatory_decision_draft_2025}.

(2.7) NEW DIMENSION (introduced in this paper): physical-AI procedure
execution choice. The patient selects which surgical robot, humanoid
companion, or rehabilitation exoskeleton performs their care from the
10 robot categories enumerated in
patients/patient_robot_instructions_fixed.tex (pages 1 through 10) and
in national-platform/patient_robot/01_preamble_robots1to5.tex
(surgical, cobot, RT positioning, needle placement, social companion)
and 02_robots6to10_closing.tex (humanoid, RT motion-tracking, imaging,
steerable needle, rehabilitation exoskeleton). Cite
\cite{kawchak_2026_18810541}, \cite{kawchak_2026_18778219} (for the
USL scoring framework that quantifies each robot's readiness for
patient choice), and \cite{target2025robotic_bronchoscopy} (for an
example of robotic-assisted bronchoscopy that demonstrates the
patient-control benefit of robot-executed procedures).

After the seven dimensions, embed a comprehensive ASCII diagram of the
seven-dimensional patient-control framework as a verbatim block. The
diagram should map each dimension to the proposed bill (HR 4501
through HR 4507) that addresses it, plus the prior law that each bill
adapts or revises. Use box-drawing characters (+, -, |, +) only and
keep every line under 80 columns to avoid right-margin overflow per
the global formatting brief in main.tex.

Close this section with one paragraph stating that for INDIVIDUAL
patients the framework defines per-decision rights (e.g., the patient
PAT-2026-0042 in patient-journey/stage_05_surgery.py selects the da
Vinci Xi for lobectomy), and for BROAD adoption across all U.S. cancer
patients in trials the framework defines population-level rights and
implementation deadlines. The bills in Sections 3 through 9 implement
both views simultaneously. Cite \cite{kawchak_2026_19119939} for the
single-patient example.]
